(1) 展开对流算符得
\[
\partial_t^2y+2U\partial_t\partial_xy+(U^2-c^2)\partial_x^2y=0.
\]
设 \(y\propto e^{i(kx-\omega t)}\)，代入得到
\[
(\omega-Uk)^2=c^2k^2.
\]
所以两条有符号频率分支为 \(\omega_\pm(k)=(U\pm c)k\)，在固定 k 的情况下可为负，不影响实场解通过正负共轭成分组成。

(2) 相速度 \(\omega_\pm/k=U\pm c\)，群速度 \(\partial\omega_\pm/\partial k=U\pm c\)，两者各自相等，表示此线性色散模型的两个波包在实验室中沿固定特征速度传播。没有其他分散尺度，所以同一分支内部不同 k 不展宽。

(3) 固定实验室正频率 \(\omega >0\)，可写
\[
k_+=\frac{\omega}{U+c},\qquad k_-=\frac{\omega}{U-c}
\]
（分母不为零时）。若 |U|<c，两个群速度 \(U+c>0,U-c<0\)，右行与左行特征分别向下游和上游传播；在 U=c 时向上游分支实验室群速度为零，不能用非零有限 k 承载给定正 \(\omega \)。若 U>c，则两支群速度都为正，即使第二支在流体静止系相对介质逆向传播，实验室仍向下游移动。

(4) 引入随流坐标 \(x'=x-Ut,t'=t\)，链式法则给 \((\partial_t+U\partial_x)=\partial_{t'}\)，方程化为
\[
\partial_{t'}^2y=c^2\partial_{x'}^2y.
\]
相位变换 \(\omega'=\omega-Uk\)，在介质系为 \(\omega'^2=c^2k^2\)，恢复通常双向波。U>c 时实验室的最慢传播特征速度仍为 \(U-c>0\)，因此上游无法收到有限传播速度的线性小信号。
